
%\documentclass[a4paper,12pt]{article}
\documentclass[10pt]{article}
\renewcommand{\baselinestretch}{1.5} 

\setlength{\parindent}{0pt}

%% Language and font encodings
\usepackage[utf8]{inputenc}
\usepackage[T1]{fontenc}
\usepackage[square,sort&compress]{natbib}    
\usepackage{booktabs}
\usepackage[french,english]{babel}
%% Sets page size and margins
\usepackage[top=1in,left=1.5in,bottom=1in,right=1in]{geometry}

\usepackage{relsize}
%% Useful packages
\usepackage{amsmath, amssymb, bm}
\usepackage{graphicx}
\usepackage{subfig}
\usepackage{authblk}
\usepackage{xcolor}


\usepackage[
  bookmarks       = true,             % Show the bookmarks
  unicode         = true,             % Use Unicode                 
  pdftoolbar      = true,             % Show Acrobat’s toolbar
  pdfmenubar      = true,             % Show Acrobat’s menu
  pdffitwindow    = false,            % Window fit to page when opened
  pdfstartview    = {FitH},           % Fit the page width to the window
  pdfnewwindow    = true,             % Open the links in a new PDF window
  colorlinks      = true,             % Use colored links
  linkcolor       = blue,             % Internal links color
  citecolor       = blue,            % Bibliographic links color
  filecolor       = cyan,             % File links color
  urlcolor        = magenta           % External links color
]{hyperref}
\usepackage{algorithm}
% \usepackage{algorithmic}
\usepackage{algorithmicx}
\usepackage{listings}
\usepackage{algcompatible}
\usepackage{algpseudocode}
\usepackage{tikz}
\usepackage{calc}
\usepackage{pgfplots}
\usepackage{pdfpages}
\pgfplotsset{compat=1.14}
\usepackage{tikzducks}
\usetikzlibrary{snakes}
\usetikzlibrary{positioning}
\usetikzlibrary{shapes,arrows}
\usetikzlibrary{automata,arrows,positioning,calc}
% \usepackage{cmbright}

\makeatletter
\newcounter{algorithmicH}% New algorithmic-like hyperref counter
\let\oldalgorithmic\algorithmic
\renewcommand{\algorithmic}{%
  \stepcounter{algorithmicH}% Step counter
  \oldalgorithmic}% Do what was always done with algorithmic environment
\renewcommand{\theHALG@line}{ALG@line.\thealgorithmicH.\arabic{ALG@line}}
\makeatother

% \geometry{inner=1cm, outer=1cm, top=1cm, bottom=0cm}
\newcommand\fillin[1]{\makebox[#1]{\dotfill}}
% Use a modern font (that is not pixelated)
\usepackage{kpfonts}
% Use microtype to improve readability
\usepackage[protrusion = true, expansion = true]{microtype} 
% Use multiple columns
\usepackage{multicol}
% \usepackage{subfigure}
% \usepackage[subfigure]{tocloft}
% Display list of Equations
% \usepackage{tocloft}
% Insert empty pages
% \usepackage{afterpage}
% Configure style aspects of the document
\usepackage{caption}            % To configure the captions
\usepackage{siunitx}            % For SI units
\usepackage{graphicx}           % To include graphics
% \usepackage[dvipsnames]{xcolor} % To configure colors
\usepackage{fancyhdr}           % To configure the header and footer
% \pagestyle{fancy}               % Use a fancy header style
% \fancyhf{}                      % Set the header format
% \fancyheadoffset{0.0 cm}        % Set the header offset
\lhead{\tfm}                    % Set the header right-side contents
\rhead{\name}                   % Set the header left-side contents
% \cfoot{\thepage}                % Set the footer page number
\usepackage{amsthm}
\usepackage[capitalise,nameinlink,noabbrev]{cleveref}
\crefname{equation}{}{}
\crefname{enumi}{}{}


\newtheorem{lemma}{Lemma}
\newtheorem{theo}{Theorem}
\newtheorem{corollary}{corollary}
\newtheorem{coro}{Corollary}
\newtheorem{prop}{Proposition}
\newtheorem{property}{Property}
\newtheorem{remark}{Remark}
\newtheorem{ex}{Example}
\newtheorem{exercise}{Exercise}
\newtheorem{definition}{Definition}
\newtheorem{hp}{Assumption}

% Distributions.
\newcommand*{\UnifDist}{\mathsf{Unif}}
\newcommand*{\ExpDist}{\mathsf{Exp}}
\newcommand*{\DepExpDist}{\mathsf{DepExp}}
\newcommand*{\GammaDist}{\mathsf{Gamma}}
\newcommand*{\LognormalDist}{\mathsf{LogNorm}}
\newcommand*{\WeibullDist}{\mathsf{Weib}}
\newcommand*{\ParetoDist}{\mathsf{Par}}
\newcommand*{\NormalDist}{\mathsf{Normal}}

\newcommand*{\GeometricDist}{\mathsf{Geom}}
\newcommand*{\NegBinomialDist}{\mathsf{NegBin}}
\newcommand*{\PoissonDist}{\mathsf{Poisson}}

% Sets of numbers.
\newcommand*{\RL}{\mathbb{R}}
\newcommand*{\NZ}{\mathbb{N}_0}
% \newcommand*{\NL}{\mathbb{N}_+}

\newcommand*{\cond}{\mid}
\newcommand*{\given}{\,;\,}

%Protocol
\newcommand*{\PoW}{PoW\@\xspace}
\newcommand*{\PoS}{PoS\@\xspace}
\newcommand*{\PoSp}{PoSp\@\xspace}
\newcommand*{\PoI}{PoI\@\xspace}
\newcommand*{\PpS}{PpS\@\xspace}
% Regarding spacing and abbreviations.
\usepackage{xspace}

% Acronyms
% \@\xspace doesn't add space if next char is punctuation
% However, these will give 2 .'s if used at end of sentence.
\newcommand*{\eg}{e.g.\@\xspace}
\newcommand*{\ie}{i.e.\@\xspace}
\newcommand*{\iid}{i.i.d.\@\xspace}
\newcommand*{\cf}{c.f.\@\xspace}
\newcommand*{\pdf}{p.d.f.\@\xspace}
\newcommand*{\pmf}{p.m.f.\@\xspace}
\newcommand*{\cdf}{c.d.f.\@\xspace}
\newcommand*{\SMC}{\textbf{SMC}\@\xspace}
\newcommand*{\MCMC}{\textbf{MCMC}\@\xspace}
\newcommand*{\Prob}{\mathbb{P}}
\newcommand*{\E}{\mathbb{E}}


\newcommand*{\iidSim}{\overset{\mathrm{i.i.d.}}{\sim}}
\newcommand*{\bt}{\bm{\theta}}
\newcommand*{\bTheta}{\bm{\Theta}}
\newcommand*{\bbeta}{\bm{\beta}}
\newcommand*{\bx}{\mathbf{x}}
\newcommand*{\bs}{\bm{s}}
\newcommand*{\bu}{\bm{u}}
\newcommand*{\bn}{\bm{n}}

% Roman versions of things.
\newcommand*{\dd}{\mathop{}\!\mathrm{d}}
\newcommand*{\e}{\mathrm{e}}
\DeclareMathOperator*{\argmax}{arg\,max}
\DeclareMathOperator*{\argmin}{arg\,min}

\newcommand*{\norm}[1]{\lVert{} #1\rVert}

% \DeclarePairedDelimiterXPP{\ind}[1]{\ind_}{\{}{\}}{}{#1}
\newcommand*{\ind}{\mathbb{I}}

\begin{document}

\title{Mining pool network and decentralization of a PoW blockchains}
\author[1]{Hansjoerg Albrecher}
\author[2]{Pierre-O. Goffard}
\author[3]{Bilel Tounsi}
\affil[1]{\footnotesize UNIL}
\affil[2]{\footnotesize Université de Strasbourg}
\affil[3]{\footnotesize UNIL}

\date{\today }

\maketitle
\vspace{3mm}

\begin{abstract}
The Proof-of-Work (PoW) blockchain relies on miners who incur constant operational costs while receiving stochastic rewards. To mitigate financial uncertainty, many miners join Pay-Per-Share (PPS) mining pools, which provide stable payouts while transferring risk to pool managers. This paper develops a two-stage analysis of miners’ strategic behavior and pool competition. In both approaches, the pool manager’s surplus is modeled as a two-sided jump process. First, using a mean-variance trade-off approach, we characterize the optimal allocation of miners across mining pools and investigate Nash equilibrium among pools. Second, we study the pool manager’s payout problem within an expected discounted dividend framework. Using techniques from actuarial risk theory, we establish and prove the optimality of a barrier strategy under specific conditions and provide numerical methods to compute the optimal barrier and associated dividends. We then analyze the corresponding Nash equilibrium in fees under this criterion. Our findings contribute to understanding how economic incentives influence miner participation and the decentralization of PoW blockchain networks.
\end{abstract}
% {\it MSC 2010:} 60G55, 60G40, 12E10. \\
% {\it Keywords:} First hitting time, blockchain, double spending problem.\\

% !TEX root = ../main.tex


\section{Introduction}\label{sec:introduction}




The Proof-of-Work (PoW) blockchain consensus mechanism, as first introduced by \citet{nakamoto2008bitcoin}, forms the foundation of decentralized cryptocurrencies such as Bitcoin. This system relies on miners who expend computational resources to validate transactions and secure the network. However, mining is characterized by a dual-risk structure: miners face a constant operational cost (electricity, hardware maintenance) while their rewards---block subsidies and transaction fees---are inherently stochastic and infrequent. This financial uncertainty poses a significant challenge, as it directly impacts the sustainability of individual miners.\\


\noindent To mitigate this volatility, miners often choose to join mining pools. These pools aggregate computational power, distribute rewards more predictably, and ensure more stable revenue streams. Among the various pooling strategies, the Pay-per-Share (PPS) model \citep{albrecher2022blockchain,schrijvers2016incentive,rosenfeld2011analysis} provides miners with a fixed payout per valid share they contribute, effectively transferring risk to the pool manager.\\


\noindent Miners must make a strategic decision regarding participation in a mining pool, balancing expected revenue against the risk of ruin. This decision can be analyzed using an expected discounted dividend approach, which evaluates the long-term viability of miners given their financial constraints. In this context, optimal dividend strategies, such as the barrier strategy \citep{avanzi2007optimal}, emerge as critical tools for sustaining miner operations and maximizing profitability. These strategies rely on results from fluctuation theory and scale functions \citep{bayraktar2013optimal} to determine optimal capital retention and payout policies.\\


\noindent While extensive research exists on individual miner strategies, the surplus process of the pool manager remains relatively unexplored. It can be modeled as a two-sided jump process, incorporating incoming miner contributions and outgoing payouts. However, existing literature provides limited results on optimal surplus management in this setting. In this paper, we establish that a band strategy is optimal in general and that a simple barrier strategy becomes under certain restrictions. We further develop numerical methods to identify the optimal barrier level and compute corresponding dividends.\\


\noindent Beyond individual miner and pool manager strategies, we leverage our findings to analyze the equilibrium distribution of miners across available mining pools. By considering a set of miners and a set of pool managers, we derive insights into how miners allocate their computational power in response to different payout schemes and risk profiles. Our study further extends to a numerical investigation of Nash equilibria, illustrating the strategic interactions between miners and pool managers in a competitive setting.\\


\noindent This research contributes to the ongoing discourse on decentralization in PoW blockchains. By providing a framework for mining pool strategies and decision-making, we offer insights into how economic incentives shape miner participation, thereby influencing the overall decentralization and stability of blockchain networks.

\input{blockchain_mining_pool}
% !TEX root = ../main.tex
\section{Wealth processes for miners and pool managers}\label{sec:mining_pool}
\subsection{miners' wealth process}\label{sec:sub_mining_pool}
As introduced by \cite{hansjoerg2022profitability}, The miners' wealth can be modeled as a stochastic process
\begin{equation}\label{eq:surplus_solo_mining}
X_t = x -c\cdot t + N_t\cdot b,\text{ for } t\geq 0,
\end{equation}
where $x>0$ represents the initial capital, $c$ is the operational cost, $b$ is the reward for discovery of blocks and $(N_t)_{t\geq0}$ is a Poisson process with intensity $\lambda$, which represents the number of blocks found per unit of time. The variance of this process exposes miners to the risk of insolvency. To reduce income volatility, miners join mining pools, aggregating their computational power to increase the frequency of rewards and obtain more stable earnings. A pool manager oversees this collective activity by assessing each miner’s contribution and allocating rewards accordingly. This paper considers the Pay-per-Share (PPS) scheme which is a widely used reward mechanism. Under this system, miners submit shares as proof of their work. A share corresponds to a solution of a cryptographic puzzle with a lower difficulty than that required to mine a full block. A PPS deal is characterized by a difficulty reduction $\delta$ $\in(0,1]$ and a pool fee $f$ $\in[0,1)$. In the PPS framework, the pool manager pays a fixed reward $b_k<b$ to miners for each submitted share. Then the stochastic process that represents the miners' wealth becomes
\begin{equation}\label{eq:surplus_solo_in_pool_mining}
\tilde{X}_t = x -c\cdot t + \tilde{N}_t\cdot b_k,\text{ for } t\geq 0,
\end{equation}
where $b_k=b\cdot\delta\cdot(1-f)$ is the share reward and $(\tilde{N}_t)_{t\geq0}$ is a Poisson process with intensity $\tilde{\lambda}=\frac{\lambda}{\delta}\geq\lambda$, which exceeds that of $(N_t)_{t\geq0}$ in (\ref{eq:surplus_solo_mining}). The advantage of the PPS system is that it transfers risk from miners to pool managers. Equations (\ref{eq:surplus_solo_mining}) and (\ref{eq:surplus_solo_in_pool_mining}) allow \citet{albrecher2022blockchain} to quantify miners’ risk through ruin probabilities. Note that solo mining can be interpreted as a particular case of the PPS framework. Indeed, it corresponds to a contract with no fee and no difficulty reduction, that is, $f = 0$ and $\delta = 1$. In this case, the share reward reduces to $b_k = b$ and the arrival intensity becomes $\tilde{\lambda} = \lambda$. Consequently, equation (\ref{eq:surplus_solo_in_pool_mining}) simplifies to
\begin{equation}
\tilde{X}_t = x - c\cdot t + \tilde{N}_t \cdot b, \quad t \geq 0,
\end{equation}
which coincides with (\ref{eq:surplus_solo_mining}). Hence, solo mining can be viewed as a particular case of PPS contract.





% JE VAIS AJOUTER çA POUR LA PARTIE DIVIDENDE ET JE VAIS COMMENCER PAR DIRE QU'IL FAUT ADAPTER LE PROCESSUS

% More generally, the previous models can be viewed as particular cases of a more general framework in which the miner allocates her computational resources across multiple pools. A miner may choose to split her computational resources across several mining pools simultaneously. Consider $n \in \mathbb{N}$ mining pools, each offering a PPS contract characterized by a difficulty reduction $\delta_k$ and a fee $f_k$, for $k = 1, \dots, n$. These contractual parameters determine both the effective share arrival rate and the associated reward. More precisely, pool $k$ induces a share arrival intensity $\lambda / \delta_k$ and a reward per share $b_k = b\,\delta_k(1 - f_k)$. In addition, the miner retains the possibility to mine solo, which can be interpreted as an artificial pool indexed by $k = 0$, with $\delta_0 = 1$, $f_0 = 0$, and thus $b_0 = b$.

% The miner allocates her total computational power $\lambda$ among the available options according to a weight vector
% \[
% w = (w_0, w_1, \dots, w_n) \in \Delta_n,
% \]
% where $\Delta_n = \left\{ w \in \mathbb{R}^{n+1} : \sum_{k=0}^n w_k = 1 \right\}$ denotes the unit simplex. The reward arrival rate associated with pool $k$ is then given by
% \[
% \mu_k = \frac{\lambda w_k}{\delta_k}, \quad k = 0, \dots, n.
% \]

% The $\mu_k$ define the intensities of conditionally independent Poisson processes given the allocation $w$. Aggregating these contributions, the total reward arrival process $(\hat{N}_t)_{t \geq 0}$ is a Poisson process with intensity
% \[
% \mu = \sum_{k=0}^n \mu_k.
% \]
% Each reward originates from pool $k$ with probability $p_k = \mu_k / \mu$ and has magnitude $b_k$. Therefore, the miner’s wealth process can be written as
% \begin{equation}
% X_t = x - c \cdot t + \sum_{i=1}^{\hat{N}_t} B_i,
% \end{equation}
% where $(B_i)_{i \geq 1}$ is a sequence of independent and identically distributed random variables satisfying
% \[
% \mathbb{P}(B = b_k) = \frac{\mu_k}{\mu} = p_k, \quad k = 0, \dots, n.
% \]

% This representation highlights that each contract affects both the arrival intensity and the reward size, which jointly determine the mean and variance of the miner’s wealth.

% \begin{figure}[htbp]
%     \centering
%     \includegraphics[width=0.7\textwidth]{multipool_diagram.png}
%     \caption{Multi-pool mining architecture.}
%     \label{fig:multipool_structure}
% \end{figure}


\subsection{Pool manager's wealth process}\label{ssec:pool_manager_process}

Consider $m \in \mathbb{N}$ mining pools, each offering a PPS contract characterized by a difficulty reduction $\delta_k$ and a fee $f_k$, for $k = 1, \dots, m$. A PPS pool $k$ attracts miners by offering a contract $(\delta_k, f_k)$.  These contractual parameters determine both the effective share arrival rate and the associated reward. More precisely, for a miner with hashrate $\lambda$, participation in pool $k$ leads to a share arrival intensity $\lambda / \delta_k$ and a reward per share $b_k = b\,\delta_k(1 - f_k)$. This contract increases the frequency of rewards while reducing their size, thereby lowering the variance of the miner’s wealth at the cost of a lower expected payoff. In addition, the miner retains the possibility to mine solo, which can be interpreted as an artificial pool indexed by $k = 0$, with $\delta_0 = 1$, $f_0 = 0$, and thus $b_0 = b$. The pool $k$ collects a total hashpower $H_k = \sum_{i \in \mathcal{I}} \lambda_i$, where $\mathcal{I}$ denotes the set of miners who joined pool $k$ and $\lambda_i$ is the hashpower contributed by miner $i$. The rate of arrival of shares of pool $k$ is given by
\[
    \mu_k = \frac{H_k}{\delta_k},
\]
which corresponds to the average number of shares found by the pool participants per time unit. Each share has probability $\delta_k/(1+\delta_k)$ of being an actual solution to the cryptographic 
puzzle, and probability $1/(1+\delta_k)$ of being a simple share. The wealth process of the pool 
manager is given by
\begin{equation}\label{eq:surplus_pool_manager}
    Y_t = y + \sum_{i=1}^{M_t} V_i,
\end{equation}
where 
\begin{itemize}
    \item $(M_t)_{t\geq 0}$ is a Poisson process with intensity $\mu_k$,
    \item the $V_i$'s are i.i.d. with probability distribution
    \[
        \mathbb{P}\!\left(V = -b_k\right) = \frac{1}{1+\delta_k}, 
        \qquad 
        \mathbb{P}\!\left(V = b - b_k\right) = \frac{\delta_k}{1+\delta_k},
    \]
    where $b_k=b(1-f_k)\delta_k$ is the reward that the pool $k$ offers.
\end{itemize}

\begin{remark}
Using the superposition, one can express the surplus of the pool manager (\ref{eq:surplus_pool_manager}) as the difference between two Poisson processes associated to deterministic positive and negative jumps.
\end{remark}


Each time a share arrives, the pool manager pays the miner the fixed PPS amount $b(1-f_k)\delta_k$, regardless of whether the share is a solution or not. When the share turns out to be an actual solution, the manager additionally collects the block reward $b$. The pool is profitable in expectation if and only if
\begin{equation}
\mathbb{E}[V] > 0,
\end{equation}
which determines admissible combinations of $(\delta,f)$. This formulation makes explicit that the PPS contract transfers risk from miners to the pool manager: increasing the frequency of rewards reduces miners' variance, while concentrating risk in the manager’s wealth process.













\section{Competition between pools with mean-variance preferences}

We consider a mining ecosystem composed of a population of heterogeneous miners and two PPS pool managers. The interaction is modeled as a static game in which for a given difficulty reduction $(\delta_1,\delta_2)$ $\in(0,1]^2$, pool managers compete through fees, anticipating the participation decisions of miners who evaluate available options according to a mean--variance criterion.

\subsection{Miners' mean-variance}

Consider a population of $n$ miners, indexed by $i = 1, \dots, n$. Each miner $i$ 
is characterized by a block discovery rate $\lambda_i > 0$, an operational cost 
$c_i \geq 0$, and a risk-aversion parameter $\gamma_i > 0$. The miner can choose between solo mining, indexed by $k=0$, and two PPS pools, indexed by $k=1,2$. For a given contractual pair $(f_k,\delta_k)$, the miner's wealth process over a time horizon $T$ follows the PPS formulation (\ref{sec:mining_pool}) introduced in Section \ref{sec:sub_mining_pool}. The miner evaluates option $k$ according to the mean--variance criterion
\begin{equation}
S_i(k)
:=
\mathbb{E}\!\left[X_T^{(i,k)}\right]
-
\gamma_i\,\mathrm{Var}\!\left(X_T^{(i,k)}\right).
\end{equation}
The expected wealth is
\begin{equation}
\mathbb{E}\!\left[X_T^{(i,k)}\right]
=
x-c_iT+T\frac{\lambda_i}{\delta_k}b(1-f_k)\delta_k
=
x-c_iT+T\lambda_i b(1-f_k),
\end{equation}
while the variance is
\begin{equation}
\mathrm{Var}\!\left(X_T^{(i,k)}\right)
=
T\frac{\lambda_i}{\delta_k}b^2(1-f_k)^2\delta_k^2
=
T\lambda_i b^2(1-f_k)^2\delta_k.
\end{equation}
Hence,
\begin{equation}
S_i(k)
=
x-c_iT
+
T\lambda_i
\left(
b(1-f_k)-\gamma_i b^2(1-f_k)^2\delta_k
\right).
\end{equation}
Since the additive term $x-cT$ and the multiplicative term $T$ are common across options, the miner's decision is independent of the horizon $T$.

For solo mining, we set $(f_0,\delta_0)=(0,1)$, which gives
\begin{equation*}
S_i(0)=\lambda_i\left(b-\gamma_i b^2\right).
\end{equation*}
For participation in pool $k\in\{1,2\}$,
\begin{equation*}
S_i(k)
=
\lambda_i\left(
b(1-f_k)-\gamma_i b^2(1-f_k)^2\delta_k
\right).
\end{equation*}
The miner therefore chooses
\begin{equation}\label{eq:miner_opti}
k_i^*\in\arg\max_{k\in\{0,1,2\}} S_i(k).
\end{equation}
Because $\lambda_i>0$ is a common multiplicative factor across all options, it does not affect the ranking of the alternatives. As expressed \citet{goffard2025hashpower}, when we use the mean-variance criterion, miner's decision depends only on $\gamma_i$ and on the contractual parameters $(f_k,\delta_k)$. Thus, miners with identical risk preferences choose the same option and allocate their entire hashrate to the pool that maximizes their objective function.

\subsection{Pool managers mean-variance}

For a given horizon $T=1$, the surplus process of pool manager $k$ associated with a contractual pair $(f_k,\delta_k)$ is the compound Poisson process (\ref{eq:surplus_pool_manager}). Its expectation is given by
\begin{equation}\label{eq:mean_pool_mana}
\mathbb{E}\!\left[Y_t^{(k)}\right]
=
y_k+\mu_k \,\frac{b\delta_k}{1+\delta_k}\big(f_k(1+\delta_k)-\delta_k\big),\\
\end{equation}
and its variance is given by
\begin{equation}\label{eq:var_pool_mana}
\mathrm{Var}\!\left(Y_t^{(k)}\right)
=
\mu_k 
\left[
\frac{b^2(1-f_k)^2\delta_k^2}{1+\delta_k}
+
\frac{b^2\delta_k\big(1-(1-f_k)\delta_k\big)^2}{1+\delta_k}
\right],
\end{equation}
where $\mu_k$ denotes the share-arrival intensity of pool $k$, as introduced in 
Section~\ref{ssec:pool_manager_process}. For a fixed difficulty reduction $\delta_k$, 
pools compete solely through their fee. Consequently, the total hashpower 
attracted by pool $k$ depends on the fee vector $(f_k, f_{-k})$, where $f_k$ is the 
fee set by pool $k$ and $f_{-k} = (f_j)_{j \neq k}$ denotes the fees of competing pools. 
We therefore write
\begin{equation}\label{eq:hashrate_f}
    H_k(f_k, f_{-k}) = \sum_{i \in \mathcal{I}_k(f_k, f_{-k})} \lambda_i,
\end{equation}
where $\mathcal{I}_k(f_k, f_{-k})$ denotes the set of miners who join pool $k$ under 
the fee profile $(f_k, f_{-k})$, and $\lambda_i$ is the hashpower contributed by miner $i$. Thus the share-arrival intensity of pool $k$ is given by

\begin{equation}\label{eq:new_pool_intensity}
    \mu_k = \frac{H_k(f_k, f_{-k})}{\delta_k}.
\end{equation}
In particular, if all miners join pool $k$, then $\mathcal{I}_k = \{1, \dots, n\}$ 
and $H_k(f_k, f_{-k}) = \lambda$, where $\lambda = \sum_{i=1}^{n} \lambda_i$ 
denotes the total hashpower available on the network. 

The manager evaluates outcomes according to a conditional mean--variance criterion. 
For a given realization of the aggregate hashrate $H_k$, the surplus process 
$Y_t^{(k)}$ is a compound Poisson process with intensity 
$\mu_k = H_k / \delta_k$. The manager's conditional objective is defined as
\begin{equation}\label{mean_var_brut_pool}
J_k(H_k)
=
\mathbb{E}\!\left[Y_t^{(k)} \mid H_k\right]
-
\eta_k\,\mathrm{Var}\!\left(Y_t^{(k)} \mid H_k\right),
\end{equation}
where $\eta_k>0$ is the manager's risk-aversion coefficient. Substituting (\ref{eq:mean_pool_mana}), (\ref{eq:var_pool_mana}) and (\ref{eq:new_pool_intensity}) into (\ref{mean_var_brut_pool}) yields
\begin{equation}
\label{eq:Jmanager_final_chapter}
J_k(H_k(f_k, f_{-k}))
=
y_k
+
H_k(f_k, f_{-k})\frac{b}{1+\delta_k}
\left[
f_k(1+\delta_k)-\delta_k
-\eta_k b\Big(
\delta_k(1-f_k)^2+\big(1-(1-f_k)\delta_k\big)^2
\Big)
\right].
\end{equation}

\subsection{Fee competition using expected approach}

We now analyze a fee-setting game between the two pool managers. The parameters $(\delta_1,\delta_2)$ are fixed ex ante and interpreted as technological characteristics of the two pools. Managers choose fees $(f_1,f_2)$, miners then respond by selecting the option that maximizes their mean--variance score \eqref{eq:miner_opti}, and managers receive payoffs given by \eqref{eq:Jmanager_final_chapter}. 

\subsubsection{Threshold structure of miner choices}

We assume that miners are heterogeneous only through their risk-aversion coefficient and that
\begin{equation}
\gamma \sim \mathrm{Unif}[0,\bar{\gamma}],
\end{equation}
independently across $n$ miners. Since $\lambda_i$ factors out of all pairwise comparisons, it is convenient to normalize the miner's score by $\lambda_i$ and define
\begin{equation}
s_0(\gamma)=b-\gamma b^2,
\end{equation}
\begin{equation}
s_k(\gamma)=b(1-f_k)-\gamma \delta_k b^2(1-f_k)^2,
\qquad k\in\{1,2\}.
\end{equation}
Each miner selects $k \in \{0,1,2\}$ to maximize $s_k(\gamma)$. Since the mean--variance scores are affine in the risk-aversion parameter $\gamma$, pairwise differences are affine, implying that optimal choices are characterized by cutoff values of $\gamma$. Consequently, the set of types selecting any given option forms an (possibly empty) interval in $[0,\bar{\gamma}]$, and the equilibrium allocation is fully described by a finite collection of such thresholds.

\begin{prop}[Threshold structure]\label{prop:threshold_structure}
Consider a miner with risk-aversion parameter $\gamma \in [0,\bar{\gamma}]$. The optimal choice problem \eqref{eq:miner_opti} is characterized by three thresholds.

\medskip

\noindent\textbf{(i) Pool versus solo mining.}  
For $k\in\{1,2\}$, pool $k$ is preferred to solo mining if and only if
\begin{equation*}
s_k(\gamma)\ge s_0(\gamma)
\quad\Longleftrightarrow\quad
\gamma \ge \gamma_{k0}(f_k),
\end{equation*}
where
\begin{equation}
\gamma_{k0}(f_k)
=
\frac{f_k}{b\left(1-\delta_k(1-f_k)^2\right)}.
\end{equation}

\medskip

\noindent\textbf{(ii) Comparison between pools.}  
Define
\begin{equation}
D_{12}(f_1,f_2)
=
\delta_1(1-f_1)^2-\delta_2(1-f_2)^2.
\end{equation}
If $D_{12}(f_1,f_2)\neq 0$, the comparison between the two pools is characterized by the threshold
\begin{equation}
\gamma_{12}(f_1,f_2)
=
\frac{f_2-f_1}{b\,D_{12}(f_1,f_2)},
\end{equation}
such that
\begin{equation}
s_1(\gamma)\ge s_2(\gamma)
\quad\Longleftrightarrow\quad
\begin{cases}
\gamma \le \gamma_{12}(f_1,f_2), & \text{if } D_{12}(f_1,f_2)>0,\\
\gamma \ge \gamma_{12}(f_1,f_2), & \text{if } D_{12}(f_1,f_2)<0.
\end{cases}
\end{equation}
If $D_{12}(f_1,f_2)=0$, the comparison between the two pools is independent of $\gamma$ and depends only on expected returns.

\medskip

\noindent\textbf{(iii) Partition of the type space.}  
The interval $[0,\bar{\gamma}]$ is partitioned into at most three subintervals, corresponding respectively to miners who optimally choose solo mining, pool~1, and pool~2.
\end{prop}

\begin{proof}
For each $k\in\{0,1,2\}$, the score $s_k(\gamma)$ is affine in $\gamma$. Hence, for any pair $(i,j)$, the difference $s_i(\gamma)-s_j(\gamma)$ is also affine and admits at most one root. Each pairwise comparison therefore induces a unique threshold value of $\gamma$, which establishes the expressions for $\gamma_{k0}$ and $\gamma_{12}$. The partition of the interval $[0,\bar{\gamma}]$ follows directly from these pairwise comparisons.
\end{proof}

\begin{remark}[Economic interpretation]\label{remark:Economic_interpretation}
The thresholds $\gamma_{k0}(f_k)$ represent entry conditions into the pooling market: only sufficiently risk-averse miners prefer joining a pool over solo mining. The threshold $\gamma_{12}(f_1,f_2)$ captures the trade-off between competing pools and reflects their relative risk profiles through the term $D_{12}(f_1,f_2)$. When $D_{12}>0$, pool~1 is more risky than pool~2 and is therefore preferred only by miners with relatively low risk aversion. Conversely, when $D_{12}<0$, pool~1 offers lower variance and attracts the most risk-averse miners. 
\end{remark}
This implies a monotone sorting of miners according to their risk preferences: the mining ecosystem endogenously segments into groups that differ only through their degree of risk aversion. Consequently, the allocation of miners across pools is fully determined by the relative positions of the thresholds $\gamma_{10}$, $\gamma_{20}$, and $\gamma_{12}$. The mechanism is illustrated in Figure~\ref{fig:gamma_thresholds}, where the intersections of the score functions determine the relevant thresholds and the induced partition of the type space.


\begin{figure}[H]
\centering
\includegraphics[width=0.75\textwidth]{gamma_thresholds.png}
\caption{Mean--variance scores as functions of the risk-aversion parameter $\gamma$. 
The intersections of the score functions determine the thresholds 
$\gamma_{10}$, $\gamma_{20}$, and $\gamma_{12}$, which partition the support 
$\gamma\in[0,\bar{\gamma}]$ into regions where miners optimally choose 
solo mining, pool 1, or pool 2.}
\label{fig:gamma_thresholds}
\end{figure}


\subsubsection{Expected participation and hashrate allocation}

The threshold structure established in Proposition~\ref{prop:threshold_structure} 
allows for a tractable characterization of miner participation. Since miners are 
heterogeneous only through their risk-aversion parameter $\gamma$, and since the 
optimal choice is determined by threshold comparisons, the allocation of miners 
across alternatives can be obtained by measuring the corresponding intervals in 
the support $[0,\bar{\gamma}]$.

We denote by $p_k(f_1,f_2)$ the probability that a randomly selected miner joins 
option $k\in\{0,1,2\}$, where $k=0$ corresponds to solo mining. Under the assumption 
that $\gamma \sim \mathrm{Unif}[0,\bar{\gamma}]$, these probabilities are proportional 
to the length of the intervals identified in Proposition~\ref{prop:threshold_structure}.

\begin{prop}[Participation probabilities]\label{prop:participation_probabilities}
Let $(f_1,f_2)\in[0,1)^2$ be a given fee profile. The participation probabilities 
are determined by the relative position of the thresholds 
$\gamma_{10}(f_1)$, $\gamma_{20}(f_2)$, and $\gamma_{12}(f_1,f_2)$.

\medskip

\noindent\textbf{(i) Case $D_{12}(f_1,f_2)>0$.}  
Pool~1 is preferred only by sufficiently weakly risk-averse miners. The participation 
probability of pool~1 is given by
\begin{equation}
p_1(f_1,f_2)
=
\frac{
\big[
\min\{\bar{\gamma},\gamma_{12}(f_1,f_2)\}
-
\max\{0,\gamma_{10}(f_1)\}
\big]_+
}{\bar{\gamma}}.
\end{equation}

\medskip

\noindent\textbf{(ii) Case $D_{12}(f_1,f_2)<0$.}  
Pool~1 is preferred by sufficiently risk-averse miners. The participation probability 
of pool~1 is given by
\begin{equation}
p_1(f_1,f_2)
=
\frac{
\big[
\bar{\gamma}
-
\max\{0,\gamma_{10}(f_1),\gamma_{12}(f_1,f_2)\}
\big]_+
}{\bar{\gamma}}.
\end{equation}

\medskip

\noindent\textbf{(iii) Case $D_{12}(f_1,f_2)=0$.}  
The comparison between the two pools is independent of $\gamma$. If $f_1\le f_2$, 
pool~1 weakly dominates pool~2, and its participation probability reduces to
\begin{equation}
p_1(f_1,f_2)
=
\frac{
\big[
\bar{\gamma}
-
\max\{0,\gamma_{10}(f_1)\}
\big]_+
}{\bar{\gamma}}.
\end{equation}
Otherwise, $p_1(f_1,f_2)=0$.

\medskip

The participation probability of pool~2 is obtained symmetrically, and the share 
of solo miners satisfies
\begin{equation}
p_0(f_1,f_2)
=
1-p_1(f_1,f_2)-p_2(f_1,f_2).
\end{equation}
\end{prop}


\begin{proof}
By Proposition~\ref{prop:threshold_structure}, each option $k\in\{0,1,2\}$ is chosen 
on an interval $I_k(f_1,f_2)\subseteq[0,\bar{\gamma}]$ defined by the corresponding 
threshold conditions. Since $\gamma \sim \mathrm{Unif}[0,\bar{\gamma}]$, we have
\[
p_k(f_1,f_2)=\mathbb{P}(\gamma\in {I}(f_1,f_2))
=\frac{\mathrm{Leb}({I}(f_1,f_2))}{\bar{\gamma}},
\]
where $\mathrm{Leb}(\cdot)$ denotes the Lebesgue measure. The expressions follow 
by identifying the relevant intervals in each case according to the ordering of 
the thresholds $\gamma_{10}$, $\gamma_{20}$, and $\gamma_{12}$.
\end{proof}

Let $n$ denote the total number of miners. The expected number of miners joining 
option $k$ is given by
\begin{equation}
\mathbb{E}[N_k(f_1,f_2)]
=
n\,p_k(f_1,f_2),
\qquad k\in\{0,1,2\}.
\end{equation}

To relate participation to aggregate mining power, we assume that individual 
hashrates $(\lambda_i)_{i=1}^n$ are independent of risk preferences and satisfy
\begin{equation}
\sum_{i=1}^n \lambda_i = \lambda,
\qquad
\lambda_i \perp \gamma_i.
\end{equation}
Under this assumption, expected hashrates are proportional to participation probabilities:
\begin{equation}
\mathbb{E}[H_k(f_1,f_2)]
=
\lambda\,p_k(f_1,f_2),
\qquad k\in\{0,1,2\}.
\end{equation}

This characterization provides a complete description of the demand side of the 
mining market. The participation probabilities are piecewise-defined functions of 
the fee profile $(f_1,f_2)$, with kinks arising when one of the thresholds 
$\gamma_{10}$, $\gamma_{20}$, or $\gamma_{12}$ reaches the boundary of the support 
$[0,\bar{\gamma}]$. As a consequence, marginal changes in fees may induce discrete 
shifts in participation through changes in the ordering of these thresholds. Moreover, the aggregation step shows that the distribution of hashrate across pools is entirely summarized by the probabilities $(p_0,p_1,p_2)$ and the total available hashrate $\lambda$. This reduction will be instrumental in the subsequent analysis of pool managers' optimal fee-setting behavior.

\subsubsection{Expected payoffs of pool managers}

Using the aggregation result $\mathbb{E}[H_k(f_1,f_2)] = \lambda\,p_k(f_1,f_2)$ established in the previous subsection, the expected payoff of pool manager $k\in\{1,2\}$ can be expressed directly as a function of the fee profile $(f_1,f_2)$. Substituting into \eqref{eq:Jmanager_final_chapter} yields
\begin{equation}\label{eq:EJ_compact_final}
\mathbb{E}[J_k(f_1,f_2)]
=
y_k
+
\lambda\,p_k(f_1,f_2)\frac{b}{1+\delta_k}B_k(f_k),
\end{equation}
where
\begin{equation}
B_k(f_k)
=
f_k(1+\delta_k)-\delta_k
-\eta_k b
\Big(
\delta_k(1-f_k)^2+(1-(1-f_k)\delta_k)^2
\Big).
\end{equation}

The representation \eqref{eq:EJ_compact_final} makes explicit that the manager’s payoff is linear in the participation probability $p_k(f_1,f_2)$ and therefore in the share of total hashrate attracted by the pool. The term $B_k(f_k)$ captures the contribution of the fee choice to the payoff per unit of hashrate, incorporating both the expected revenue and the variance penalty induced by the mean--variance criterion.

This formulation highlights the trade-off faced by pool managers. Increasing the fee $f_k$ raises the expected revenue component of $B_k(f_k)$, but simultaneously reduces the participation probability $p_k(f_1,f_2)$ by making the pool less attractive to miners. In addition, the quadratic terms in $B_k(f_k)$ reflect the impact of risk preferences, penalizing fee levels that increase the variability of the surplus process.

Consequently, the competition between pools can be interpreted as a pricing game with endogenous demand: each manager chooses a fee to balance the margin extracted from participating miners against the total hashrate attracted. This reduced-form representation provides the basis for the characterization of the Nash equilibrium in fees developed in the next subsection.

\subsubsection{Nash equilibrium in a fee competition}

We now formalize the strategic interaction between pool managers. The competitive setting considered in this paper is a static two-stage game. In the first stage, each pool manager $k\in\{1,2\}$ simultaneously chooses a fee $f_k \in [0,1)$. In the second stage, miners observe the fee profile $(f_1,f_2)$ and select the option that maximizes their mean--variance criterion. This induces participation probabilities $p_k(f_1,f_2)$ and determines the allocation of hashrate across pools.

The resulting interaction can therefore be interpreted as a pricing game with endogenous demand, in which each manager anticipates the miners’ response when selecting fees.

\begin{definition}[Nash equilibrium]
A fee profile $(f_1^*,f_2^*) \in [0,1)^2$ is a Nash equilibrium if, for each $k \in \{1,2\}$, the fee $f_k^*$ solves
\begin{equation}
\sup_{f_k \in [0,1)} \; \mathbb{E}[J_k(f_k,f_{-k}^*)],
\end{equation}
where $f_{-k}^*$ denotes the equilibrium fee of the competing pool. Equivalently, a Nash equilibrium is a fee pair such that no pool manager can increase its expected payoff by unilaterally deviating from $f_k^*$, given the fee chosen by the other pool.
\end{definition}

The main difficulty in characterizing such an equilibrium lies in the fact that the participation probabilities $p_k(f_1,f_2)$ are piecewise-defined functions of the fee profile. Their analytical expressions depend on the ordering of the thresholds $\gamma_{10}$, $\gamma_{20}$, and $\gamma_{12}$, as well as on the sign of $D_{12}(f_1,f_2)$. As a consequence, the expected payoff functions are generally not globally smooth, and the analysis must be conducted region by region.

We focus on an interior coexistence region in which both pools attract a strictly positive fraction of miners. Consider the region defined by
\begin{equation}\label{eq:region_conditions_refined}
D_{12}(f_1,f_2)<0,
\qquad
0<\gamma_{12}(f_1,f_2)<\bar{\gamma},
\qquad
\gamma_{12}(f_1,f_2)\geq \gamma_{10}(f_1),
\qquad
\gamma_{12}(f_1,f_2)\geq \gamma_{20}(f_2).
\end{equation}
In this configuration, the participation probabilities simplify to
\begin{equation}
p_1(f_1,f_2)
=
\frac{\bar{\gamma}-\gamma_{12}(f_1,f_2)}{\bar{\gamma}},
\qquad
p_2(f_1,f_2)
=
\frac{\gamma_{12}(f_1,f_2)-\gamma_{20}(f_2)}{\bar{\gamma}}.
\end{equation}

Within this region, the expected payoff functions are differentiable, and interior equilibria can be characterized using first-order conditions. For manager $k$, the optimality condition is given by
\begin{equation}\label{eq:foc_general_refined}
\frac{\partial \mathbb{E}[J_k]}{\partial f_k}(f_1,f_2)
=
\lambda\frac{b}{1+\delta_k}
\left[
\frac{\partial p_k}{\partial f_k}(f_1,f_2)\,B_k(f_k)
+
p_k(f_1,f_2)\,B_k'(f_k)
\right]
=0.
\end{equation}

This condition admits a natural economic interpretation. The first term captures the marginal effect of the fee on participation (demand effect), while the second term captures the marginal effect on per-unit profitability (margin effect). At equilibrium, these two forces exactly offset each other. 
Using the expressions of $p_k(f_1,f_2)$ in the coexistence region, one obtains explicit formulas for the derivatives $\partial p_k/\partial f_k$ and $B_k'(f_k)$ (see Appendix \ref{appendix:derivatives}), and hence a system of two coupled equations

\begin{equation}
\begin{aligned}
\frac{\partial \mathbb{E}[J_1]}{\partial f_1}(f_1^*,f_2^*) &= 0, \\
\frac{\partial \mathbb{E}[J_2]}{\partial f_2}(f_1^*,f_2^*) &= 0.
\end{aligned}
\end{equation}


A candidate equilibrium $(f_1^*,f_2^*)$ must therefore solve this system. However, this is not sufficient. Since the expressions for $p_1$ and $p_2$ are valid only within the region defined by \eqref{eq:region_conditions_refined}, any solution must also satisfy the consistency conditions
\begin{equation}
D_{12}(f_1^*,f_2^*)<0,
\qquad
0<\gamma_{12}(f_1^*,f_2^*)<\bar{\gamma},
\qquad
\gamma_{12}(f_1^*,f_2^*)\geq \gamma_{10}(f_1^*),
\qquad
\gamma_{12}(f_1^*,f_2^*)\geq \gamma_{20}(f_2^*).
\end{equation}

These constraints ensure that the candidate solution is consistent with the underlying participation structure. If any of these conditions fails, the equilibrium does not belong to the interior coexistence region, and the analysis must be carried out in the corresponding alternative region of the fee space.
\begin{ex}\label{ex:1}
We consider a fee-setting game with two PPS pools and a continuum of miners whose risk-aversion parameter is uniformly distributed on $[0,\bar{\gamma}]$. The parameters are set as
\[
b=3.125,
\qquad
\bar{\gamma}=0.4,
\qquad
\lambda=100,
\]
\[
\delta_1=0.008,
\qquad
\delta_2=0.08,
\qquad
\eta_1=0.01,
\qquad
\eta_2=0.001,
\qquad
y_1=y_2=100.
\]

Solving the first-order conditions in the interior coexistence region yields the candidate fee profile
\[
f_1^*=9.9949901714\%,
\qquad
f_2^*=8.5458873160\%,
\]

and the solution is consistent with the interior coexistence region. The corresponding allocation of hashrate across solo mining and the two pools is illustrated in Figure~\ref{fig:hashrate_pie}, which reports the proportion of total hashrate collected by each option.

\begin{figure}[H]
\centering
\includegraphics[width=0.4\textwidth]{hashrate_pie.png}
\caption{Distribution of the hashrate power among solo, pool 1 and pool2.}
\label{fig:hashrate_pie}
\end{figure}
\end{ex}


The equilibrium in Example \ref{ex:1} therefore reflects the endogenous segmentation of miners according to their risk preferences. This example also highlights a limitation of the mean--variance framework in the present setting. Since each miner selects the single alternative that maximizes its objective function, the allocation rule is of all-or-nothing type: a miner can join one pool or mine solo, but cannot diversify its hashrate across several pools. From an economic perspective, this restriction is strong, as diversification is a natural response to risk. This limitation motivates the introduction of the dividend-based approach in the next chapter, which provides a richer framework for the analysis of mining decisions and pool competition.

\subsection{Sensitivity Analysis Around the Baseline Calibration}

We now complement the baseline equilibrium computation with a local sensitivity analysis around the reference calibration. The objective is not to establish global existence over the whole parameter space, but rather to assess the numerical stability and economic robustness of the equilibrium in a neighborhood of the benchmark values.

Our baseline specification is given by
\[
b=3.125,\qquad \bar{\gamma}=0.4,\qquad \Lambda=100,\qquad y_1=y_2=100,
\]
together with
\[
\delta_1=0.008,\qquad \delta_2=0.08,\qquad \eta_1=0.01,\qquad \eta_2=0.001.
\]
On the competition domain \(f_1,f_2\in[0,0.15]\), the validated equilibrium is
\[
f_1^*\approx 0.09995,\qquad f_2^*\approx 0.08546,
\]
with participation shares
\[
p_0\approx 0.07327,\qquad p_1\approx 0.80816,\qquad p_2\approx 0.11857,
\]
and expected manager payoffs
\[
E[J_1]\approx 115.47,\qquad E[J_2]\approx 100.32.
\]

\paragraph{Sensitivity with respect to \(\eta_1\) and \(\eta_2\).}

We first study how the equilibrium varies when the managers' risk-aversion parameters change locally around the baseline configuration. Since the low-fee equilibrium of the initial example is not present over the entire \((\eta_1,\eta_2)\)-space, we restrict attention to a local rectangle on which the equilibrium remains numerically validated throughout the grid. More precisely, we consider
\[
\eta_1 \in [0.008,\,0.02],\qquad \eta_2 \in [0.001,\,0.005].
\]

Over this region, we compute the surfaces associated with the equilibrium objects
\[
E[J_1](\eta_1,\eta_2),\quad f_1^*(\eta_1,\eta_2),\quad p_1^*(\eta_1,\eta_2),
\]
and
\[
E[J_2](\eta_1,\eta_2),\quad f_2^*(\eta_1,\eta_2),\quad p_2^*(\eta_1,\eta_2).
\]
The numerical evidence shows that the equilibrium varies smoothly over this neighborhood, with no breakdown of validation on the displayed grid. Hence, the baseline outcome is locally robust with respect to perturbations in the managers' risk-aversion parameters.

From an economic viewpoint, the main qualitative asymmetry of the benchmark is preserved throughout this region: pool~1 retains the larger participation share and the larger expected payoff, while pool~2 remains active but comparatively smaller. In other words, local changes in \(\eta_1\) and \(\eta_2\) modify the levels of payoffs, fees, and participation rates, but do not overturn the basic ranking implied by the reference calibration.

\paragraph{Sensitivity with respect to \(\delta_1\) and \(\delta_2\).}

We next analyze the local dependence of the equilibrium on the pool-specific variance parameters \(\delta_1\) and \(\delta_2\), while keeping the risk-aversion parameters fixed at their baseline values,
\[
\eta_1=0.01,\qquad \eta_2=0.001.
\]
Again, the analysis is conducted on a local rectangle chosen so that the equilibrium remains validated throughout the grid. Specifically, we consider
\[
\delta_1 \in [0.006,\,0.012],\qquad \delta_2 \in [0.065,\,0.095].
\]

On this domain, we compute the equilibrium surfaces
\[
E[J_1](\delta_1,\delta_2),\quad f_1^*(\delta_1,\delta_2),\quad p_1^*(\delta_1,\delta_2),
\]
and
\[
E[J_2](\delta_1,\delta_2),\quad f_2^*(\delta_1,\delta_2),\quad p_2^*(\delta_1,\delta_2).
\]
These figures provide a local characterization of how the equilibrium responds to changes in the pools' variance-reduction technologies. As in the \((\eta_1,\eta_2)\)-analysis, the computed surfaces remain regular over the displayed region, which indicates that the benchmark equilibrium is also locally robust with respect to moderate perturbations in \(\delta_1\) and \(\delta_2\).

This second exercise is economically useful because the parameters \(\delta_k\) directly govern the mean-variance trade-off offered by each pool. Varying \(\delta_1\) and \(\delta_2\) therefore changes both the attractiveness of each pool to miners and the managers' profit margins. The numerical results confirm that these effects can be substantial in level, while still preserving the existence of a well-defined low-fee equilibrium in a neighborhood of the baseline calibration.

\paragraph{Interpretation.}

Taken together, these local sensitivity exercises show that the benchmark equilibrium of the initial example is not an isolated numerical artifact. Instead, it persists over non-trivial neighborhoods in both the \((\eta_1,\eta_2)\) and \((\delta_1,\delta_2)\) dimensions. This reinforces the interpretation of the baseline outcome as a structurally meaningful equilibrium configuration of the model, rather than a knife-edge point in parameter space.

At the same time, the analysis should be interpreted as local. The displayed surfaces do not claim that a validated equilibrium exists for all admissible parameter values. Rather, they show that, around the benchmark calibration, the equilibrium is stable enough to support comparative-statics conclusions on expected payoffs, equilibrium fees, and participation shares.





















% \section{Competition between pools with dividend strategy}

\subsection{Wealth processes with diversification}
\subsubsection{miners' wealth process}

Consider $m \in \mathbb{N}$ mining pools, each offering a PPS contract characterized by a difficulty reduction $\delta_k$ and a fee $f_k$, for $k = 1, \dots, m$. Consider also a population of $n$ miners, indexed by $i = 1, \dots, n$. The wealth process \eqref{eq:surplus_solo_in_pool_mining} can be viewed as particular cases of a more general framework in which the miner allocates his computational resources $\lambda_i$ across multiple pools ($k = 1,\cdots,m$) and solo ($k=0$) simultaneously, according to a weight vector
\[
w = (w_{i,0}, w_{i,1}, \dots, w_{i,m}) \in \Delta_{i,m},
\]
where $\Delta_{i,m} = \left\{ w \in \mathbb{R_+}^{m+1} : \sum_{k=0}^m w_{i,k} = 1 \right\}$ denotes the unit simplex. The reward arrival rate of miner $i$ associated with pool $k$ is then given by
\[
\lambda_{i,k} = \frac{\lambda_i w_{i,k}}{\delta_k}, \quad k = 0, \dots, m.
\]

The $\lambda_{i,k}$ define the intensities of conditionally independent Poisson processes given the allocation $w$. Aggregating these contributions, the total reward arrival process of the miner $(\hat{N}^{(i)}_t)_{t \geq 0}$ is a Poisson process with intensity
\[
\hat{\lambda}_i = \sum_{k=0}^m \lambda_{i,k}.
\]
Each reward originates from pool $k$ with probability $p_{i,k} = \lambda_{i,k} / \hat{\lambda} _i$ and has magnitude $b_k = b(1-f_k)\delta_k$. Therefore, the miner’s wealth process can be written as
\begin{equation}\label{welath_multi_miner}
X_t = x - c \cdot t + \sum_{j=1}^{\hat{N}^{(i)}_t} B^{(i)}_j,
\end{equation}
where $(B_j)_{j \geq 1}$ is a sequence of independent and identically distributed random variables satisfying
\[
\mathbb{P}(B^{(i)} = b_k) = \frac{\lambda_{i,k}} {\hat{\lambda}_i} = p_{i,k}, \quad k = 0, \dots, m.
\]

Figure \ref{fig:multipool_structure} illustrates how the wealth process evolves in this architecture.

\begin{figure}[htbp]
    \centering
    \includegraphics[width=0.7\textwidth]{multipool_diagram.png}
    \caption{Multi-pool mining architecture.}
    \label{fig:multipool_structure}
\end{figure}

\subsubsection{Pool managers' wealth process}

Under this allocation, the total hashpower received by pool $k$ is given by
\[
H_k(w)=\sum_{i=1}^n \lambda_i w_{i,k}.
\]
As in the single-pool participation setting, the arrival rate of shares in pool $k$ is determined by the effective difficulty parameter $\delta_k$, yielding a Poisson intensity
\[
\mu_k(w)=\frac{H_k(w)}{\delta_k}
=\frac{1}{\delta_k}\sum_{i=1}^n \lambda_i w_{i,k}.
\]

Each share submitted to pool $k$ is either a valid block with probability $\delta_k/(1+\delta_k)$, or a non-block share with probability $1/(1+\delta_k)$. Regardless of its nature, the pool manager pays a fixed PPS reward $b_k=b\,\delta_k(1-f_k)$ to the contributing miner. In the event of a valid block, the manager additionally receives the block reward $b$. The wealth process of the manager of pool $k$ is therefore given by the compound Poisson process
\begin{equation}\label{eq:surplus_pool_manager_split}
Y_t^{(k)} = y_k + \sum_{j=1}^{M_t^{(k)}} V_j^{(k)},
\end{equation}
where, conditionally on the allocation profile $w$,
\begin{itemize}
    \item $(M_t^{(k)})_{t\geq 0}$ is a Poisson process with intensity $\mu_k(w)$,
    \item the sequence $(V_j^{(k)})_{j\geq 1}$ is i.i.d. with distribution
    \[
    \mathbb{P}\!\left(V^{(k)}=-b_k\right)=\frac{1}{1+\delta_k},
    \qquad
    \mathbb{P}\!\left(V^{(k)}=b-b_k\right)=\frac{\delta_k}{1+\delta_k}.
    \]
\end{itemize}

We emphasize that, in contrast to the single-pool participation case, the total hashpower $H_k(w)$ is no longer determined by a partition of miners, but rather by the aggregation of fractional contributions. This induces a coupling across pools through the allocation profile $w$, while preserving conditional independence of the underlying Poisson processes.



\subsection{Expected discounted dividend maximization}
\subsubsection{miners' dividend strategy}

A natural performance criterion for the miner is given by the classical dividend optimization problem from risk theory. The wealth process \eqref{welath_multi_miner} can be interpreted as the surplus of the miner, from which cash flows are extracted over time in the form of dividends. Formally, a dividend strategy is described by a non-decreasing, adapted process $(L_t)_{t \geq 0}$, representing the cumulative payments made up to time $t$. The controlled surplus process then writes
\[
    U_t = X_t - L_t,
\]
and ruin occurs at the stopping time
\[
    \tau = \inf\{t \geq 0 : U_t < 0\}.
\]
The objective of the miner is to maximize the expected discounted dividends
until ruin, namely
\begin{equation}\label{eq:dividend_value}
    V(x) = \sup_{L} \mathbb{E}_x\left[\int_0^{\tau} e^{-rt} \, dL_t \right],
\end{equation}
where $r \geq 0$ is a discount rate. For any fixed allocation profile $w \in \Delta_{i,m}$, the process $(X_t)_{t \geq 0}$ is a spectrally positive L\'{e}vy process. In this setting, it is well established that the optimal dividend strategy is of barrier type (see \citet{avanzi2007optimal} and \citet{bayraktar2013optimal}): there exists a level $a_w^* > 0$ such that it is optimal to pay out immediately any surplus above $a_w^*$, while no dividends are paid when the surplus lies below this level. The controlled process is thus reflected at the barrier $a_w^*$.

Denoting by $V_w(x)$ the value function associated with \eqref{eq:dividend_value} under the optimal barrier strategy, the miner's problem reduces to a static optimization over the allocation set, since each $w \in \Delta_{i,m}$ uniquely determines an optimal barrier $a_w^*$ and a corresponding value $V_w(x)$. The optimal allocation is therefore
\[
    w^* = \arg\max_{w \in \Delta_{i,m}} V_w(x).
\]
The contribution of \citet{goffard2025hashpower} is to show that this problem can be tractably analyzed in the context of PPS mining pools. By exploiting the L\'{e}vy structure of the wealth process, the value function $V_w(x)$ can be characterized in terms of scale functions. This representation allows for the computation of both the optimal barrier level $a_w^*$ and the associated value $V_w(x)$, thereby enabling a systematic comparison of PPS contracts and a characterization of the optimal allocation $w^*$.

\subsubsection{pool managers' dividend strategy}

Having characterized the wealth process \eqref{eq:surplus_pool_manager_split} of the pool
manager, we now turn to the associated dividend optimization problem. The surplus process of
pool $k$ is a pure jump process with two-sided increments: upward jumps of size $b - b_k$ occur
upon block discovery, while downward jumps of size $b_k$ occur upon submission of a non-block
share. Since the process is constant between consecutive jump times, no strategic decision needs
to be made in continuous time. It is therefore sufficient to inspect the process at jump times only. Let $(T_i^{(k)})_{i \geq 1}$ denote the jump times of $(M_t^{(k)})_{t \geq 0}$, which form a
sequence of i.i.d.\ exponential random variables with parameter $\mu_k(w)$. The embedded
discrete-time process
\[
    S_n^{(k)} := Y_{T_n^{(k)}}^{(k)} = y_k + \sum_{i=1}^n V_i^{(k)}, \quad n \geq 0,
\]
is a standard random walk. A dividend strategy is then a sequence of decisions
$\mathcal{S} = \{f_1, f_2, \ldots\}$, where each $f_n : \mathbb{R} \to \mathbb{R}_+$ specifies
the dividend paid at step $n$, leading to the controlled process
\[
    R_n^{(k)} = S_{n-1}^{(k)} - f_{n-1}\!\left(S_{n-1}^{(k)}\right) + V_n^{(k)},
    \quad n \geq 1.
\]
Ruin can only occur at jump times, and is defined as
\[
    \tau^{(k)} = \inf\!\left\{n \geq 1 : R_n^{(k)} < 0\right\}.
\]
The objective of the pool manager is to maximize the expected discounted dividends
\[
    V^{(k)}(y_k) = \mathbb{E}_{y_k}\!\left[
        \sum_{n=1}^{\tau^{(k)}} e^{-r T_n^{(k)}} f_n\!\left(S_{n-1}^{(k)}\right)
    \right],
\]
where $r \geq 0$ is the discount rate. This framework coincides with that of
\citet{albrecher2011randomized}, in which the optimal strategy is shown to be a
\textit{band strategy}: dividends are paid according to a state-dependent rule determined by a
finite collection of intervals. Band strategies are, however, difficult to handle analytically. To obtain a more tractable structure, we impose the following simplifying assumptions, which are natural in the context of
a homogeneous PPS pool.

\begin{hp}\label{ass:integer_state}
The state space of $(S_n^{(k)})_{n \geq 0}$ and $(R_n^{(k)})_{n \geq 0}$ is $\mathbb{Z}$.
\end{hp}

Assumption~\ref{ass:integer_state} requires that $b_k \in \mathbb{N}$, that the initial surplus
$y_k$ is proportional to $b_k$, and that $b$ is proportional to $b_k$. Under this assumption, $(S_n^{(k)})_{n \geq 0}$ is an integer-valued random walk with increments taking
values in $\{-1,\, b-b_k\}$ using the taxonomy of \citet{gerber1972games} for which a barrier strategy is optimal as shown in \citet{miyasawa1961economic}: there
exists a threshold $a_k^* \in \mathbb{N}$ such that the pool manager pays out all surplus above
$a_k^*$ as dividends, while retaining the full surplus when it lies below this level. We now characterize explicitly the value function and the optimal barrier level.

Define the vector of monomials
\begin{equation}\label{eq:rho_vector}
    \bm{\rho}^{[y]} := \left(\rho_1^y,\, \rho_2^y,\, \dots,\, \rho_k^y\right)^\top \in \mathbb{C}^k,
\end{equation}
where $\rho_1, \dots, \rho_k \in \mathbb{C}$ denote the $k$ roots of the characteristic equation
\begin{equation}\label{eq:characteristic}
    (r + \mu)\,\rho \;-\; \mu p\,\rho^{k} \;-\; \mu q \;=\; 0,
\end{equation}
with $p = \delta/(1+\delta)$, $q = 1/(1+\delta)$, and $\mu = \mu_k(w)$ the Poisson intensity
defined in \eqref{eq:surplus_pool_manager_split}.

\begin{prop}\label{prop:optimal_barrier}
Let $(\widetilde Y_t^{(k)})_{t \geq 0}$ denote the normalized surplus process of pool $k$ defined by
\[
\widetilde Y_t^{(k)} := \frac{Y_t^{(k)}}{b_k}.
\]
Under Assumption~\ref{ass:integer_state}, it is an integer-valued compound Poisson process with increments in $\{-1,\, k-1\}$, where $k = b/b_k \in \mathbb{N}_{\geq 2}$. Assume
moreover that the drift condition
\begin{equation}\label{eq:drift_condition}
    p(k-1) - q > 0
\end{equation}
holds, ensuring that the process is not absorbed at ruin almost surely. For a fixed
candidate barrier $a \in \mathbb{N}$ with $a \geq k - 1$, the value function under
the barrier strategy at level $a$ takes the form
\begin{equation}\label{eq:value_function_form}
    V^{(k)}(y,\, a) \;=\; \bm{\rho}^{[y]} \cdot \mathbf{c}(a), \quad y = 0, 1, \dots, a,
\end{equation}
where the coefficient vector $\mathbf{c}(a) \in \mathbb{C}^k$ is the unique solution of the
linear system
\begin{equation}\label{eq:linear_system}
    \mathbf{A}(a)\,\mathbf{c}(a) \;=\; \mathbf{d}(a).
\end{equation}
The matrix $\mathbf{A}(a) \in \mathbb{C}^{k \times k}$ and the vector
$\mathbf{d}(a) \in \mathbb{C}^{k}$ are defined row-wise by
\begin{equation}\label{eq:A_matrix}
    \mathbf{A}(a)_j \;=\;
    \begin{cases}
        (r+\mu)\,\bm{\rho}^{[y_j]} - \mu p\,\bm{\rho}^{[y_j + k - 1]},
        & j = 1, \\[0.6em]
        (r+\mu)\,\bm{\rho}^{[y_j]} - \mu p\,\bm{\rho}^{[a]}
        - \mu q\,\bm{\rho}^{[y_j - 1]},
        & j = 2, \dots, k,
    \end{cases}
\end{equation}
\begin{equation}\label{eq:d_vector}
    \mathbf{d}(a)_j \;=\;
    \begin{cases}
        0,
        & j = 1, \\[0.6em]
        \mu p\,(y_j + k - 1 - a),
        & j = 2, \dots, k,
    \end{cases}
\end{equation}
where $y_1 = 0$ and $y_j = a - (k-1) + (j-1)$ for $j = 2, \dots, k$.
The optimal barrier is then given by
\begin{equation}\label{eq:optimal_barrier}
    a_k^* \;=\; \operatorname*{arg\,max}_{a \,\in\, \mathbb{N},\; a \,\geq\, k-1}
    V^{(k)}(y,\, a).
\end{equation}
\end{prop}

\begin{proof}
We derive the value function by analyzing the dynamic programming equation satisfied
by $V^{(k)}(\cdot, a)$ over an infinitesimal time interval $dt$. Since the surplus
process is constant between jump times, the only events to consider are the arrival
or non-arrival of a jump during $[0, dt)$.

\medskip
\noindent\textbf{Step 1: Dynamic programming equation.}
With probability $1 - \mu\,dt + o(dt)$, no jump occurs and the value function is
discounted at rate $r$. With probability $\mu p\,dt + o(dt)$, an upward jump of
size $k-1$ occurs, and with probability $\mu q\,dt + o(dt)$, a downward jump of
size $1$ occurs. Letting $dt \to 0$ and rearranging, one obtains the
difference equation
\begin{equation}\label{eq:diff_eq}
    (r + \mu)\,V^{(k)}(y,a)
    \;=\;
    \mu p\,V^{(k)}(y + k - 1,\, a)
    \;+\;
    \mu q\,V^{(k)}(y - 1,\, a),
    \quad y = 1, \dots, a-(k-1),
\end{equation}
which holds in the interior of the state space, away from the boundaries.
This is a linear homogeneous difference equation of order $k$, whose general
solution is of the form
\begin{equation}\label{eq:general_solution}
    V^{(k)}(y, a) \;=\; \sum_{i=1}^{k} c_i(a)\,\rho_i^{\,y}
    \;=\; \bm{\rho}^{[y]} \cdot \mathbf{c}(a),
\end{equation}
where $\rho_1, \dots, \rho_k$ are the roots of the characteristic
equation~\eqref{eq:characteristic} and $\mathbf{c}(a) = (c_1(a), \dots, c_k(a))^\top$
are coefficients to be determined by boundary conditions (see e.g. \citet{jerri2013linear}).

\medskip
\noindent\textbf{Step 2: Boundary conditions at $y = 0$.}
For $y = 0$, a downward jump leads to immediate ruin, so that
$V^{(k)}(-1, a) = 0$. The difference equation~\eqref{eq:diff_eq} then
simplifies to
\begin{equation}\label{eq:bc_ruin}
    (r + \mu)\,V^{(k)}(0,\, a) \;=\; \mu p\,V^{(k)}(k-1,\, a).
\end{equation}
Substituting the general solution~\eqref{eq:general_solution} into~\eqref{eq:bc_ruin}
yields the first row of the linear system, namely
\begin{equation*}
    \left((r+\mu)\,\bm{\rho}^{[0]} - \mu p\,\bm{\rho}^{[k-1]}\right)
    \cdot \mathbf{c}(a) \;=\; 0,
\end{equation*}
which corresponds to row $j = 1$ of~\eqref{eq:A_matrix}--\eqref{eq:d_vector}
with $y_1 = 0$.

\medskip
\noindent\textbf{Step 3: Boundary conditions near the barrier.}
For $y \in \{a-(k-1)+1, \dots, a\}$, an upward jump of size $k-1$ brings the
surplus above the barrier. Under the barrier strategy, the excess is immediately
paid as a lump-sum dividend and the surplus is reset to $a$, so that
\begin{equation*}
    V^{(k)}(y + k - 1,\, a)
    \;\longrightarrow\;
    V^{(k)}(a,\, a) + (y + k - 1 - a).
\end{equation*}
Substituting into~\eqref{eq:diff_eq} gives, for each such $y$,
\begin{equation}\label{eq:bc_barrier}
    (r+\mu)\,V^{(k)}(y,\,a)
    \;=\;
    \mu p\!\left(V^{(k)}(a,\,a) + y + k - 1 - a\right)
    + \mu q\,V^{(k)}(y-1,\,a).
\end{equation}
Substituting the general solution~\eqref{eq:general_solution}
into~\eqref{eq:bc_barrier} and rearranging yields, for $j = 2, \dots, k$,
\begin{equation*}
    \left(
        (r+\mu)\,\bm{\rho}^{[y_j]}
        - \mu p\,\bm{\rho}^{[a]}
        - \mu q\,\bm{\rho}^{[y_j-1]}
    \right) \cdot \mathbf{c}(a)
    \;=\;
    \mu p\,(y_j + k - 1 - a),
\end{equation*}
with $y_j = a-(k-1)+(j-1)$, which corresponds to rows $j = 2, \dots, k$
of~\eqref{eq:A_matrix}--\eqref{eq:d_vector}.

\medskip
\noindent\textbf{Step 4: Determination of the optimal barrier.}
Combining Steps 2 and 3 yields the complete $k \times k$ linear
system~\eqref{eq:linear_system}, from which the coefficient vector $\mathbf{c}(a)$
is uniquely determined for each candidate barrier $a \geq k-1$.
The corresponding value function $V^{(k)}(y, a)$ is then evaluated via
\eqref{eq:value_function_form}. The optimal barrier $a_k^*$ is defined
as in~\eqref{eq:optimal_barrier}, and is obtained numerically by solving
the system over a finite grid of candidate values and selecting the maximizer.
\end{proof}

\begin{ex}\label{ex:barrier}
We illustrate numerically the characterization of the optimal dividend barrier
for a PPS pool manager in the lattice setting of
Proposition~\ref{prop:optimal_barrier}.

The parameters are set as
\[
b = 40,
\qquad
f = 0.4,
\qquad
k = 10,
\]
\[
\mu = 6,
\qquad
r = 0.05,
\qquad
x_0 = 20.
\]

The difficulty parameter is chosen so that the surplus process evolves on the
lattice $b_k \mathbb{Z}$, namely
\[
\delta = \frac{1}{k(1-f)},
\]
which implies $b_k = b/k = 4$ and normalized increments in $\{-1,\,k-1\}$.

For each candidate barrier level $a$, the value function $V(x_0,a)$ is computed
analytically by solving the linear system associated with the boundary conditions
of Proposition~\ref{prop:optimal_barrier}. The optimal barrier is obtained by
maximizing $V(x_0,a)$ over a finite grid, yielding
\[
a^* = 108,
\qquad
V(x_0,a^*) \approx 77.3.
\]

To validate this result, we estimate the value function via Monte Carlo simulation of the controlled surplus process under the barrier strategy. The jump times are generated according to a Poisson process of intensity $\mu$, and discounted dividends are accumulated until ruin. Confidence intervals at the $95\%$ level are computed for each barrier. Figure~\ref{fig:barrier_validation} compares the analytical value function with the Monte Carlo estimates.

\begin{figure}[H]
\centering
\includegraphics[width=0.6\textwidth]{div_analytic.png}
\caption{Analytical value function $V(x_0,a)$ and Monte Carlo 95\% confidence intervals.}
\label{fig:barrier_validation}
\end{figure}
\end{ex}



\input{blockchain_equilibrium}



\section*{Acknowledgmements} 
\bibliography{ref}
\bibliographystyle{plainnat}

\appendix
\appendix

\section{Derivatives of Participation Functions and Payoff Components}\label{appendix:derivatives}

\subsection{Interior coexistence region}

We consider the interior coexistence region characterized by
\[
D_{12}(f_1,f_2)<0,
\qquad
0<\gamma_{12}(f_1,f_2)<\bar{\gamma},
\qquad
\gamma_{12}(f_1,f_2)\ge \gamma_{10}(f_1),
\qquad
\gamma_{20}(f_2)\le \gamma_{12}(f_1,f_2).
\]

In this region, participation probabilities are given by
\begin{equation}
p_1(f_1,f_2)=\frac{\bar{\gamma}-\gamma_{12}(f_1,f_2)}{\bar{\gamma}},
\qquad
p_2(f_1,f_2)=\frac{\gamma_{12}(f_1,f_2)-\gamma_{20}(f_2)}{\bar{\gamma}}.
\end{equation}

\subsection{Derivative of $\gamma_{12}$}

Recall that
\begin{equation}
\gamma_{12}(f_1,f_2)
=
\frac{f_2-f_1}{bD_{12}(f_1,f_2)},
\qquad
D_{12}(f_1,f_2)=\delta_1(1-f_1)^2-\delta_2(1-f_2)^2.
\end{equation}

Differentiating with respect to $f_1$ yields
\begin{equation}
\frac{\partial \gamma_{12}}{\partial f_1}
=
\frac{
-\,bD_{12}(f_1,f_2)
+
2\delta_1(1-f_1)b(f_2-f_1)
}{
b^2D_{12}(f_1,f_2)^2
}.
\end{equation}

Differentiating with respect to $f_2$ yields
\begin{equation}
\frac{\partial \gamma_{12}}{\partial f_2}
=
\frac{
bD_{12}(f_1,f_2)
-
2\delta_2(1-f_2)b(f_2-f_1)
}{
b^2D_{12}(f_1,f_2)^2
}.
\end{equation}

\subsection{Derivative of $\gamma_{20}$}

Recall that
\begin{equation}
\gamma_{20}(f_2)
=
\frac{f_2}{b\left(1-\delta_2(1-f_2)^2\right)}.
\end{equation}

Differentiating with respect to $f_2$ yields
\begin{equation}
\frac{\partial \gamma_{20}}{\partial f_2}
=
\frac{
1-\delta_2(1-f_2)^2
-
2\delta_2 f_2(1-f_2)
}{
b\left(1-\delta_2(1-f_2)^2\right)^2
}.
\end{equation}

\subsection{Derivatives of participation probabilities}

Using the previous results, we obtain
\begin{equation}
\frac{\partial p_1}{\partial f_1}
=
-\frac{1}{\bar{\gamma}}
\frac{\partial \gamma_{12}}{\partial f_1}
=
\frac{
D_{12}(f_1,f_2)
-
2\delta_1(1-f_1)(f_2-f_1)
}{
b\bar{\gamma}\,D_{12}(f_1,f_2)^2
}.
\end{equation}

Similarly,
\begin{equation}
\frac{\partial p_2}{\partial f_2}
=
\frac{1}{\bar{\gamma}}
\left(
\frac{\partial \gamma_{12}}{\partial f_2}
-
\frac{\partial \gamma_{20}}{\partial f_2}
\right),
\end{equation}
that is,
\begin{align}
\frac{\partial p_2}{\partial f_2}
&=
\frac{1}{\bar{\gamma}}
\left[
\frac{
D_{12}(f_1,f_2)
-
2\delta_2(1-f_2)(f_2-f_1)
}{
bD_{12}(f_1,f_2)^2
}
-
\frac{
1-\delta_2(1-f_2)^2
-
2\delta_2 f_2(1-f_2)
}{
b\left(1-\delta_2(1-f_2)^2\right)^2
}
\right].
\end{align}

\subsection{Derivatives of $B_k(f_k)$}

Recall that
\begin{equation}
B_k(f_k)
=
f_k(1+\delta_k)-\delta_k
-\eta_k b
\Big(
\delta_k(1-f_k)^2+(1-(1-f_k)\delta_k)^2
\Big).
\end{equation}

Differentiating with respect to $f_k$ yields
\begin{equation}
B_k'(f_k)
=
(1+\delta_k)
-
2\eta_k b\delta_k
\Big((1+\delta_k)f_k-\delta_k\Big).
\end{equation}



\end{document}

